% GENERATED FILE. Edit canonical lessons, then run npm run generate:books.
% canonical-source-hash: fnv1a64:beb3aeaeaa6af7f3
% canonical-lessons: TE-C374-vyakhyanam, TE-C374-vyakhyata, TE-C374-srota, TE-C374-tilakincu, TE-C374-mukhamukhi

\chapter{On Air}
\label{ch:te-on-air}
\hlchaptermodality{374}

\begin{chapteropening}
\textbf{By the end of this chapter:} \emph{I can say a commentary, a commentator, a listener, to behold and a face-to-face interview.}

Complete the last lesson of chapter 374: I can say a commentary, a commentator, a listener, to behold and a face-to-face interview.
\end{chapteropening}

\section[\texorpdfstring{\te{వ్యాఖ్యానం} (vyākhyānaṁ)}{వ్యాఖ్యానం (vyākhyānaṁ)}]{\te{వ్యాఖ్యానం} (vyākhyānaṁ) — a commentary}

\label{lesson:TE-C374-vyakhyanam}



\begin{quote}
Before the new one: say the Telugu for a burrakatha ballad, then the Telugu for to make a start.
\end{quote}

\subsection*{You'll want to know: \te{వ్యాఖ్యానం}}
\textbf{\te{వ్యాఖ్యానం}} — \emph{vyākhyānaṁ} — \textquotedblleft{}a commentary\textquotedblright{}.

\begin{quote}
\te{రేడియోలో} \te{క్రికెట్} \te{వ్యాఖ్యానం} \te{వినిపిస్తోంది}.
\end{quote}

\begin{quote}
\emph{rēḍiyōlō krikeṭ vyākhyānaṁ vinipistōndi.}
\end{quote}

\begin{quote}
Cricket commentary is playing on the radio.
\end{quote}

\begin{cousinweb}
\textbf{\te{వ్యాఖ్యానం}} (\emph{vyākhyānaṁ}) is borrowed from Sanskrit \emph{vyākhyāna}, \textquotedblleft{}comment, narration\textquotedblright{} (CDIAL 12188).
\end{cousinweb}

\begin{grammarlens}[title={What you've built}]
One more word for telling what happened.
\end{grammarlens}

\subsection*{Your turn}
\begin{itemize}
\raggedright
  \item \emph{Say it:} \emph{vyākhyānaṁ}
  \item \emph{Say it:} \emph{vyākhyānaṁ}, once more
  \item \emph{Say it:} \emph{rēḍiyōlō krikeṭ vyākhyānaṁ vinipistōndi.} — Cricket commentary is playing on the radio.
  \item \emph{Recall:} say the Telugu for a burrakatha ballad, then the Telugu for to make a start, then say \emph{vyākhyānaṁ} again
\end{itemize}

\begin{tcolorbox}[breakable,colback=teal!4,colframe=teal!35!black,title={Before you move on}]
What does \te{వ్యాఖ్యానం} mean? (\textquotedblleft{}A commentary\textquotedblright{}.) Say it once more.
\end{tcolorbox}
\section[\texorpdfstring{\te{వ్యాఖ్యాత} (vyākhyāta)}{వ్యాఖ్యాత (vyākhyāta)}]{\te{వ్యాఖ్యాత} (vyākhyāta) — a commentator; a presenter}

\label{lesson:TE-C374-vyakhyata}



\begin{quote}
Before the new one: say the Telugu for to make a start, then the Telugu for a commentary.
\end{quote}

\subsection*{You'll want to know: \te{వ్యాఖ్యాత}}
\textbf{\te{వ్యాఖ్యాత}} — \emph{vyākhyāta} — \textquotedblleft{}a commentator; a presenter\textquotedblright{}.

\begin{quote}
\te{ఆ} \te{వ్యాఖ్యాత} \te{మాటలు} \te{అందరికీ} \te{నచ్చాయి}.
\end{quote}

\begin{quote}
\emph{ā vyākhyāta māṭalu andarikī naccāyi.}
\end{quote}

\begin{quote}
Everyone liked that commentator's words.
\end{quote}

Plural: \te{వ్యాఖ్యాతలు}, \emph{vyākhyātalu}.

The person behind a \te{వ్యాఖ్యానం}, a commentary, on radio or TV.

\begin{cousinweb}
\textbf{\te{వ్యాఖ్యాత}} (\emph{vyākhyāta}) is borrowed from Sanskrit \emph{vyākhyātṛ}, \textquotedblleft{}an explainer, a commentator\textquotedblright{} (Monier-Williams), from the same root as \te{వ్యాఖ్యానం}.
\end{cousinweb}

\begin{grammarlens}[title={What you've built}]
One more word for telling what happened.
\end{grammarlens}

\subsection*{Your turn}
\begin{itemize}
\raggedright
  \item \emph{Say it:} \emph{vyākhyāta}
  \item \emph{Say it:} \emph{vyākhyāta}, once more
  \item \emph{Say it:} \emph{ā vyākhyāta māṭalu andarikī naccāyi.} — Everyone liked that commentator's words.
  \item \emph{Recall:} say the Telugu for to make a start, then the Telugu for a commentary, then say \emph{vyākhyāta} again
\end{itemize}

\begin{tcolorbox}[breakable,colback=teal!4,colframe=teal!35!black,title={Before you move on}]
What does \te{వ్యాఖ్యాత} mean? (\textquotedblleft{}A commentator; a presenter\textquotedblright{}.) Say it once more.
\end{tcolorbox}
\section[\texorpdfstring{\te{శ్రోత} (śrōta)}{శ్రోత (śrōta)}]{\te{శ్రోత} (śrōta) — a listener}

\label{lesson:TE-C374-srota}



\begin{quote}
Before the new one: say the Telugu for a commentary, then the Telugu for a commentator.
\end{quote}

\subsection*{You'll want to know: \te{శ్రోత}}
\textbf{\te{శ్రోత}} — \emph{śrōta} — \textquotedblleft{}a listener\textquotedblright{}.

\begin{quote}
\te{రేడియో} \te{శ్రోతలు} \te{చాలా} \te{ఉత్తరాలు} \te{పంపారు}.
\end{quote}

\begin{quote}
\emph{rēḍiyō śrōtalu cālā uttarālu pampāru.}
\end{quote}

\begin{quote}
Radio listeners sent many letters.
\end{quote}

Plural: \te{శ్రోతలు}, \emph{śrōtalu}.

\begin{cousinweb}
\textbf{\te{శ్రోత}} (\emph{śrōta}) is borrowed from Sanskrit \emph{śrotṛ}, \textquotedblleft{}one who hears, a hearer\textquotedblright{} (Monier-Williams).
\end{cousinweb}

\begin{grammarlens}[title={What you've built}]
One more word for telling what happened.
\end{grammarlens}

\subsection*{Your turn}
\begin{itemize}
\raggedright
  \item \emph{Say it:} \emph{śrōta}
  \item \emph{Say it:} \emph{śrōta}, once more
  \item \emph{Say it:} \emph{rēḍiyō śrōtalu cālā uttarālu pampāru.} — Radio listeners sent many letters.
  \item \emph{Recall:} say the Telugu for a commentary, then the Telugu for a commentator, then say \emph{śrōta} again
\end{itemize}

\begin{tcolorbox}[breakable,colback=teal!4,colframe=teal!35!black,title={Before you move on}]
What does \te{శ్రోత} mean? (\textquotedblleft{}A listener\textquotedblright{}.) Say it once more.
\end{tcolorbox}
\section[\texorpdfstring{\te{తిలకించు} (tilakiñcu)}{తిలకించు (tilakiñcu)}]{\te{తిలకించు} (tilakiñcu) — to behold; to watch}

\label{lesson:TE-C374-tilakincu}



\begin{quote}
Before the new one: say the Telugu for a commentator, then the Telugu for a listener.
\end{quote}

\subsection*{You'll want to know: \te{తిలకించు}}
\textbf{\te{తిలకించు}} — \emph{tilakiñcu} — \textquotedblleft{}to behold; to watch\textquotedblright{}.

\begin{quote}
\te{వేలాది} \te{మంది} \te{ఆ} \te{ఉత్సవాన్ని} \te{తిలకించారు}.
\end{quote}

\begin{quote}
\emph{vēlādi mandi ā utsavānni tilakiñcāru.}
\end{quote}

\begin{quote}
Thousands watched that celebration.
\end{quote}

Past (I …): \te{తిలకించాను}, \emph{tilakiñcānu}.

A news word for watching an event or a show.

\begin{cousinweb}
\textbf{\te{తిలకించు}} (\emph{tilakiñcu}) is a written-style Telugu verb for watching a spectacle; its origin is not settled.
\end{cousinweb}

\begin{grammarlens}[title={What you've built}]
One more word for telling what happened.
\end{grammarlens}

\subsection*{Your turn}
\begin{itemize}
\raggedright
  \item \emph{Say it:} \emph{tilakiñcu}
  \item \emph{Say it:} \emph{tilakiñcu}, once more
  \item \emph{Say it:} \emph{vēlādi mandi ā utsavānni tilakiñcāru.} — Thousands watched that celebration.
  \item \emph{Recall:} say the Telugu for a commentator, then the Telugu for a listener, then say \emph{tilakiñcu} again
\end{itemize}

\begin{tcolorbox}[breakable,colback=teal!4,colframe=teal!35!black,title={Before you move on}]
What does \te{తిలకించు} mean? (\textquotedblleft{}To behold; to watch\textquotedblright{}.) Say it once more.
\end{tcolorbox}
\section[\texorpdfstring{\te{ముఖాముఖి} (mukhāmukhi)}{ముఖాముఖి (mukhāmukhi)}]{\te{ముఖాముఖి} (mukhāmukhi) — a face-to-face interview}

\label{lesson:TE-C374-mukhamukhi}



\begin{quote}
Before the new one: say the Telugu for a listener, then the Telugu for to behold.
\end{quote}

\subsection*{You'll want to know: \te{ముఖాముఖి}}
\textbf{\te{ముఖాముఖి}} — \emph{mukhāmukhi} — \textquotedblleft{}a face-to-face interview\textquotedblright{}.

\begin{quote}
\te{టీవీలో} \te{ఆ} \te{నటుడితో} \te{ముఖాముఖి} \te{వచ్చింది}.
\end{quote}

\begin{quote}
\emph{ṭīvīlō ā naṭuḍitō mukhāmukhi vaccindi.}
\end{quote}

\begin{quote}
A face-to-face interview with that actor was on TV.
\end{quote}

Literally face to face; \te{ఇంటర్వ్యూ} says the same.

\begin{cousinweb}
\textbf{\te{ముఖాముఖి}} (\emph{mukhāmukhi}) is built on Sanskrit \emph{mukha}, \textquotedblleft{}the mouth, the face\textquotedblright{} (CDIAL 10158), said twice: face to face.
\end{cousinweb}

\begin{grammarlens}[title={What you've built}]
One more word for telling what happened.
\end{grammarlens}

\subsection*{Your turn}
\begin{itemize}
\raggedright
  \item \emph{Say it:} \emph{mukhāmukhi}
  \item \emph{Say it:} \emph{mukhāmukhi}, once more
  \item \emph{Say it:} \emph{ṭīvīlō ā naṭuḍitō mukhāmukhi vaccindi.} — A face-to-face interview with that actor was on TV.
  \item \emph{Recall:} say the Telugu for a listener, then the Telugu for to behold, then say \emph{mukhāmukhi} again
\end{itemize}

\begin{tcolorbox}[breakable,colback=teal!4,colframe=teal!35!black,title={Before you move on}]
What does \te{ముఖాముఖి} mean? (\textquotedblleft{}A face-to-face interview\textquotedblright{}.) Say it once more.
\end{tcolorbox}
